%% Version 3.1 December 2024
%% Springer Nature journal article template conversion for Engineering with Computers
%% Prepared from the Engineering with Computers final manuscript.

\documentclass[lineno,pdflatex,sn-mathphys-num]{sn-jnl}

%%%% Standard packages
\usepackage{graphicx}%%
\usepackage{multirow}%%
\usepackage{amsmath,amssymb,amsfonts}%%
\usepackage{bm}%%
\usepackage{booktabs,longtable,array}%%
\usepackage{xcolor}%%
\usepackage{colortbl}%%
\usepackage{textcomp}%%
\usepackage{manyfoot}%%
\usepackage{algorithm}%%
\usepackage{algorithmicx}%%
\usepackage{algpseudocode}%%
\usepackage{listings}%%
\usepackage{subcaption}%%
\usepackage[capitalise,noabbrev]{cleveref}%%

\graphicspath{{./}}
\definecolor{highlightbg}{RGB}{255,250,205}

\raggedbottom

\begin{document}

\begin{titlepage}
\thispagestyle{empty}
\small
\begin{center}
{\Large\bfseries Title Page}\\[0.7em]
{Engineering with Computers -- Original Article}\\[1.0em]
{\LARGE\bfseries A GA-Warm-Started Mixed-Integer Nonlinear Framework for Discrete IS~800-Compliant Space-Truss Optimization}\\[0.5em]
{\footnotesize Short title: A GA-warm-started MINLP framework for code-constrained truss optimization}
\end{center}

\vspace{0.7em}
\noindent\textbf{Authors:} First Author\textsuperscript{1,*} and Second Author\textsuperscript{1}\\
\textbf{Affiliation:} \textsuperscript{1}Department of Civil Engineering, Institute/University Name, City, State, India.\\
\textbf{Corresponding author:} First Author; \texttt{corresponding.author@example.com}\\
\textbf{Contributing author:} Second Author; \texttt{second.author@example.com}\\
\textbf{ORCID IDs:} First Author: [ORCID ID]; Second Author: [ORCID ID]

\vspace{0.7em}
\noindent\textbf{Acknowledgements.} The authors gratefully acknowledge the computational support provided by the Department of Civil Engineering.

\vspace{0.45em}
\noindent\textbf{Funding.} The authors received no specific funding for this work. Update this statement if any institutional, government, or project funding supported the study.

\vspace{0.45em}
\noindent\textbf{Author contributions.} Author 1: Conceptualisation, methodology, software, validation, formal analysis, visualization, and writing--original draft. Author 2: Supervision, methodology review, validation review, and writing--review and editing.

\vspace{0.45em}
\noindent\textbf{Competing interests.} The authors declare that they have no known competing financial interests or personal relationships that could have appeared to influence the work reported in this paper.

\vspace{0.45em}
\noindent\textbf{Ethics approval and consent to participate.} Not applicable.

\vspace{0.45em}
\noindent\textbf{Consent for publication.} Not applicable.

\vspace{0.45em}
\noindent\textbf{Data availability.} The benchmark definitions, section catalogue data, optimization results, and post-processing files required to reproduce the tables and figures will be made available in a public repository at: [repository DOI/URL to be added before submission].

\vspace{0.45em}
\noindent\textbf{Code availability.} The Python implementation used for the direct-stiffness analysis, IS~800:2007 member checks, optimization algorithms, statistical runs, and figure/table generation will be made available in the same repository at: [repository DOI/URL to be added before submission].

\vspace{0.45em}
\noindent\textbf{Use of AI-assisted technologies.} During the preparation of this manuscript, the authors used Claude (Anthropic) and ChatGPT (OpenAI) to assist with \LaTeX{} formatting, language editing, manuscript organisation, and journal-specific wording. The authors reviewed and edited all AI-assisted content and take full responsibility for the scientific accuracy and integrity of the work.

\vfill
\noindent\textit{Author action before submission: replace all placeholder names, emails, affiliation, ORCID IDs, and repository DOI/URL.}
\end{titlepage}
\clearpage
\setcounter{page}{1}

\begin{center}
{\LARGE\bfseries A GA-Warm-Started Mixed-Integer Nonlinear Framework for Discrete IS~800-Compliant Space-Truss Optimization}
\end{center}

\section*{Abstract}
Discrete sizing and topology optimization of steel space trusses becomes a
non-smooth mixed-integer problem when catalogue sections and member buckling
checks are imposed during the search. This study develops a
GA-warm-started mixed-integer nonlinear optimization framework for Indian
standard equal-angle trusses. A genetic algorithm first explores binary
member topology and SP~6(1):1964 section indices. Elite feasible or
near-feasible designs then warm-start a second integer nonlinear search in
which catalogue indices are refined while IS~800:2007 stress, slenderness,
and Annex~D compression-capacity checks are evaluated explicitly for each
candidate design. Structural response is computed using the direct stiffness
method, so each candidate is assessed through analysis--design coupling
rather than post-optimization section rounding. The framework is tested on
6-, 25-, and 72-bar space-truss benchmarks and compared with deterministic,
mathematical-programming, and metaheuristic baselines over 20 independent
stochastic runs. The proposed hybrid reproduces the 6-bar benchmark best
weight of 101.17~kg, reduces the 25-bar GA best baseline by 31.8\% to 248.10~kg,
and reduces the 72-bar GA best baseline by 34.6\% to 241.67~kg. The results show
that warm-started integer nonlinear search can improve catalogue-constrained
truss designs while maintaining direct code-based feasibility screening. The
study also identifies cases where lighter penalty-based designs require
independent post-optimization verification, emphasizing the importance of
transparent feasibility reporting in code-constrained structural
optimization.

\vspace{0.75em}
\noindent\textbf{Keywords:} Space truss optimization, Mixed-integer nonlinear optimization, Genetic algorithm, Integer differential evolution, IS~800:2007, Discrete steel sections

\vspace{1em}
\section*{Article Highlights}
\begin{itemize}
  \item GA warm-starting improves catalogue-based truss optimization.
  \item IS~800 buckling checks are evaluated during the search.
  \item Integer section indices avoid post-optimization rounding.
\end{itemize}

\section{Introduction}
\label{sec:intro}
%% ══════════════════════════════════════════════════════════════════

Space trusses are three-dimensional pin-jointed structural systems that
carry external actions primarily through axial member forces. Their low
self-weight and high stiffness-to-mass ratio make them attractive for
long-span roofs, towers, industrial frames, and modular spatial structures.
From a computational viewpoint, their design is a useful benchmark class
because a small change in member area or member existence can alter the
stiffness matrix, force redistribution, buckling demand, and feasibility of
other members. The resulting optimization problem is therefore non-convex,
mixed-discrete, and strongly coupled through the structural analysis model.

Classical truss optimization separates the problem into sizing, topology,
and shape variables. In the present benchmark formulation, nodal geometry is
kept fixed so that the computational contribution can be isolated to two
coupled decisions: (i) binary member topology, which decides whether a
candidate member remains active; and (ii) catalogue-discrete sizing, which
selects a steel section for each active member. This narrower formulation is
important for transparent benchmarking, because it avoids claiming shape or
geometry optimization while still retaining the main mixed-integer challenge
encountered in catalogue-based steel design.

In the Indian steel-design context, member sizing must satisfy
IS~800:2007~\citep{IS800}, and practical sections are selected from
SP~6(1):1964~\citep{SP61}. These requirements make the search space
fundamentally discrete: a continuous optimal area cannot simply be rounded
without changing stiffness, stress, slenderness, and buckling capacity.
Furthermore, compression members must be checked using the IS~800 Annex~D
column-buckling formulation rather than a single constant allowable stress.
Thus, a computational workflow intended for code-constrained truss design
should evaluate member feasibility during the optimization loop and should
report the final demand-to-capacity ratios, not only a penalized objective
value.

\subsection{Review of existing methods}

\subsubsection{Classical mathematical-programming and resizing methods}

Early structural-optimization work established the foundations of automated
minimum-weight truss design through linearized analysis and member resizing.
The optimal-structure formulation of \citet{Dorn1964} and the fully stressed
design concepts summarized by \citet{Gallagher1973} remain useful because
they are fast and physically interpretable. Sequential approximation methods,
including the approximation concepts of \citet{Schmit1974} and the mixed-variable
dual approach of \citet{Fleury1986}, made nonlinear structural synthesis more
tractable, while discrete structural-optimization formulations such as
\citet{Ringertz1988} introduced direct treatment of catalogue-like choices.
These methods are attractive for rapid screening, but their linearizations or
continuous relaxations can lose feasibility after catalogue rounding, especially
when member-force redistribution is strong.

Mixed-integer linear and nonlinear programming provide a more formal route to
discrete section selection. However, truss problems with nonlinear stress,
slenderness, and buckling constraints are difficult because the feasible set is
non-smooth after catalogue indexing and topology decisions. General MINLP
reviews emphasize the high computational cost of such problems and the need for
problem-specific reformulations or good initializations~\citep{Belotti2013}.
For this reason, the present paper uses the term MINLP to describe the problem
class and uses a warm-started integer nonlinear search as the practical Phase~2
solver, rather than claiming a deterministic global MINLP solution.

\subsubsection{Metaheuristic and population-based truss optimization}

Metaheuristic methods have become common in truss optimization because they can
work directly with integer variables and non-differentiable penalties. Simulated
annealing~\citep{Kirkpatrick1983}, genetic algorithms~\citep{Holland1975},
particle swarm optimization~\citep{Kennedy1995}, ant-colony optimization
~\citep{Dorigo1997}, differential evolution~\citep{Storn1997}, and
Big Bang--Big Crunch optimization~\citep{Erol2006} are representative examples
of search strategies that can be adapted to discrete structural design.
Within space-truss benchmarking, \citet{Camp2004} used ant-colony optimization
for 10-bar and 25-bar examples, \citet{Camp2007} used Big Bang--Big Crunch
optimization for space trusses, and \citet{Perez2007} demonstrated a particle
swarm approach for structural design optimization. \citet{Shallan2014} further
used genetic algorithms with innovative variables for space-truss optimization.
These studies provide useful benchmark baselines, but most of them use AISC or
generic constraint settings rather than Indian steel-code catalogue checks.

Penalty functions are convenient in metaheuristics, but they also create an
important reporting issue. A design with the best penalized fitness is not
necessarily the same as the best design after an independent member-by-member
code check. A very large penalty coefficient can reduce this risk, but it does
not remove the need to report the final feasibility status. This distinction is
central to the present paper: the proposed workflow is not claimed to dominate
all penalty methods in weight, but it makes the feasibility-screening protocol
explicit and repeatable for IS~800/SP~6(1) benchmark studies.

\subsubsection{Recent computational-topology context}


Recent computational structural-optimization literature has increasingly emphasized
reproducible algorithmic workflows, benchmark transparency, and fair comparison
between optimization strategies. In this context, the present study focuses on a
narrower but practically important problem: catalogue-discrete steel space-truss
optimization with explicit IS~800:2007 feasibility screening. This does not claim to
replace broader topology-optimization or lattice-design methods; rather, it addresses
the specific need for a reproducible code-constrained steel-design workflow in which
national code checks are embedded directly in candidate evaluation.

Scientific-computing reproducibility is also increasingly important for such
studies. The present implementation is based on open Python scientific-computing
tools, including SciPy's optimization and numerical routines~\citep{Virtanen2020}
and the HiGHS linear-programming interface associated with modern simplex and
interior-point solvers~\citep{Huangfu2018}. This implementation choice supports
replicable benchmark generation, but the final submission must still provide a
code/data repository and a complete parameter table so that reviewers can rerun
the examples.

\subsection{Research gap}
\label{sec:gap}

The literature therefore leaves a specific and defensible gap. Existing
space-truss optimization studies provide strong benchmarks for metaheuristic
sizing and topology search, and recent computational-topology studies provide
advanced methods for lattice and continuum-like design. However, there is still
limited reproducible evidence for a workflow that combines all of the following
within one benchmark package:
\begin{enumerate}
  \item binary member-topology decisions for candidate space-truss members;
  \item native SP~6(1):1964 integer section indices without post-optimization
        area rounding;
  \item direct IS~800:2007 stress, slenderness, and Annex~D compression checks
        during candidate evaluation; and
  \item statistical comparison against deterministic, mathematical-programming,
        and metaheuristic baselines under the same analysis model.
\end{enumerate}
This gap is deliberately narrower than claiming the first truss-optimization
method or the first mixed-integer topology optimizer. The novelty lies in the
Indian-code, catalogue-discrete, feasibility-traceable computational workflow
and in the GA warm-start used to improve the subsequent integer nonlinear
search.

\subsection{Contributions of this paper}

This paper makes four contributions aligned with a computational-engineering
journal scope:
\begin{enumerate}
  \item It formulates a two-phase GA-warm-started mixed-integer nonlinear
        framework for topology and sizing optimization of fixed-geometry space
        trusses with discrete Indian standard equal-angle sections.
  \item It embeds explicit IS~800:2007 Annex~D compression-capacity evaluation
        and slenderness checks in the candidate-evaluation workflow, so that
        final designs are reported only after direct member-level feasibility
        screening.
  \item It provides a reproducible benchmark comparison of 14 deterministic,
        mathematical-programming, and metaheuristic methods on 6-, 25-, and
        72-bar space trusses over 20 independent stochastic runs.
  \item It reports where penalty-based methods can produce lighter candidates
        that require independent post-optimization checking, thereby separating
        weight performance from feasibility-reporting guarantees.
\end{enumerate}

The remainder of this paper is organised as follows.
\Cref{sec:problem} states the optimization problem formally.
\Cref{sec:methods} describes the baseline methods and the proposed
GA--MINLP workflow.
\Cref{sec:benchmarks} defines the benchmark trusses.
\Cref{sec:results} presents statistical results over 20 independent runs.
\Cref{sec:discussion} discusses the findings.
\Cref{sec:conclusion} concludes.


%% ══════════════════════════════════════════════════════════════════
\section{Problem Formulation}
\label{sec:problem}
%% ══════════════════════════════════════════════════════════════════

\subsection{General optimisation statement}

The space truss weight minimisation problem is:
\begin{equation}
  \min_{\mathbf{x}}\; W(\mathbf{x})
  = \sum_{i=1}^{n_m} \rho\,A_i(\mathbf{x})\,L_i(\mathbf{x})
  \label{eq:objective}
\end{equation}
subject to:
\begin{align}
  \sigma_i &\leq f_{\mathrm{allow},i}(\mathbf{x}),
    \quad \forall\,i = 1,\ldots,n_m
    \tag{stress} \label{eq:stress} \\
  \frac{KL_i}{r_i} &\leq
    \begin{cases}180 & \text{(compression)}\\400 & \text{(tension)}\end{cases}
    \tag{slenderness / IS~800 Cl.~7.1} \label{eq:slend} \\
  \delta_{\max} &\leq \delta_{\mathrm{allow}}
    \tag{deflection} \label{eq:defl} \\
  A_i &\in \mathcal{A}_{\mathrm{cat}} \subset \mathrm{SP\;6(1):1964}
    \tag{discrete catalogue} \label{eq:disc} \\
  y_i &\in \{0,1\},\quad \forall\,i
    \tag{topology} \label{eq:topo}
\end{align}
where $\rho = 7850~\mathrm{kg/m^3}$ is the steel density,
$f_y = 250~\mathrm{MPa}$, $\gamma_{m0} = 1.10$,
$\mathcal{A}_{\mathrm{cat}}$ is the set of 51 SP~6(1) cross-sectional
areas, and $y_i=1$ if member $i$ is present.

\subsection{IS~800:2007 Annex~D: design compressive stress}

The design compressive stress $f_{cd}$ for a member with effective
slenderness $KL/r$ is computed by the Annex~D buckling curve as:
\begin{align}
  f_{cc} &= \frac{\pi^2 E}{(KL/r)^2}
    \label{eq:fcc} \\
  \bar{\lambda} &= \sqrt{\frac{f_y}{f_{cc}}}
    \label{eq:lambda} \\
  \phi &= 0.5\!\left[1 + \alpha\!\left(\bar{\lambda}-0.2\right)
          + \bar{\lambda}^2\right]
    \label{eq:phi} \\
  f_{cd} &= \frac{f_y}{\gamma_{m0}
    \!\left[\phi + \sqrt{\phi^2 - \bar{\lambda}^2}\right]}
    \label{eq:fcd}
\end{align}
Here $\alpha=0.49$ for buckling class~c angle sections.
For tension members, $f_{\mathrm{allow,t}} = f_y/\gamma_{m0}$.
The demand-to-capacity ratio is defined as:
\begin{equation}
  \mathrm{DCR}_i = \frac{\sigma_i}{f_{\mathrm{allow},i}} \leq 1.0
  \label{eq:dcr}
\end{equation}

\subsection{Structural analysis: Direct Stiffness Method}

Member forces are obtained by the Direct Stiffness Method~(DSM)
\citep{McGuire2000}.
The local stiffness matrix for member $i$ of area $A_i$, length $L_i$,
and direction cosines $[\ell,\,m,\,n]$ is:
\begin{equation}
  k_i = \frac{EA_i}{L_i}
  \begin{bmatrix}
    \ell^2 & \ell m & \ell n & -\ell^2 & -\ell m & -\ell n \\
    & m^2 & mn & -\ell m & -m^2 & -mn \\
    & & n^2 & -\ell n & -mn & -n^2 \\
    & & & \ell^2 & \ell m & \ell n \\
    & & & & m^2 & mn \\
    \text{sym} & & & & & n^2
  \end{bmatrix}
  \label{eq:klocal}
\end{equation}
The global stiffness matrix $\mathbf{K}$ is assembled by standard
direct stiffness assembly.
Boundary conditions are applied by elimination; free-DOF displacements
are obtained by solving $\mathbf{K}_{ff}\mathbf{U}_f = \mathbf{F}_f$.


%% ══════════════════════════════════════════════════════════════════
\section{Optimization Methods}
\label{sec:methods}
%% ══════════════════════════════════════════════════════════════════

\subsection{Classical deterministic methods}

\subsubsection{Fully Stressed Design (FSD)}
The FSD heuristic iteratively rescales each member area to its
stress-optimal value and selects the nearest SP~6(1) section:
\begin{equation}
  A_i^{(t+1)} = A_i^{(t)}\cdot
    \frac{\sigma_i^{(t)}}{f_{\mathrm{allow},i}}
  \label{eq:fsd}
\end{equation}
Slenderness is enforced by bumping the section index upward until
$KL_i/r_i$ satisfies \cref{eq:slend}.
FSD converges in 10--30 iterations but handles sizing only and
converges to a locally optimal design.

\subsubsection{Linear Programming (LP)}
The linearised problem minimises $\sum_i \rho L_i A_i$ subject to
$A_i \geq |F_i|/f_{\mathrm{allow},i}$, solved by the HiGHS simplex
solver~\citep{Huangfu2018}.
Iterative re-solves with updated forces (Sequential LP, SLP) capture
stiffness redistribution.

\subsubsection{Sequential Quadratic Programming (SQP)}
At each iterate $\mathbf{x}_k$, the BFGS Hessian approximation $H_k$
defines a quadratic subproblem:
\begin{equation}
  \min_{\mathbf{d}}\;\tfrac{1}{2}\mathbf{d}^T H_k\mathbf{d}
    + \nabla f(\mathbf{x}_k)^T\mathbf{d}
  \quad\text{s.t.}\;\nabla g_i(\mathbf{x}_k)^T\mathbf{d}
    + g_i(\mathbf{x}_k)\leq 0
  \label{eq:sqp}
\end{equation}
Implemented via \texttt{scipy.optimize.minimize(method='SLSQP')}
\citep{Virtanen2020}.

\subsubsection{Mixed-Integer Linear Programming (MILP)}
Binary selection variables $y_{g,k}\in\{0,1\}$ for section $k$ in
member group $g$:
\begin{equation}
  \min \sum_{g=1}^{n_g}\sum_{k=1}^{n_k} w_k L_g y_{g,k}
  \quad\text{s.t.}\;
  \sum_k y_{g,k}=1,\;
  \sum_k A_k y_{g,k}\geq \frac{|F_g|}{f_{\mathrm{allow}}}
  \label{eq:milp}
\end{equation}
Solved by HiGHS branch-and-bound \citep{Huangfu2018}.

\subsubsection{Mixed-Integer Nonlinear Programming (MINLP)}
The complete discrete IS~800 problem:
\begin{equation}
  \min_{\mathbf{k}}\;W(\mathbf{k})
  \quad\text{s.t.}\;g_i(\mathbf{k})\leq 0\;\forall i,\;
  k_j\in\{0,\ldots,50\}
  \label{eq:minlp}
\end{equation}
where $g_i$ includes the exact Annex~D buckling constraint.
Solved by Differential Evolution with \texttt{integrality=1}
\citep{Virtanen2020}.

\subsection{Metaheuristic methods}

\subsubsection{Simulated Annealing (SA)}
Acceptance probability and geometric cooling schedule:
\begin{equation}
  P(\text{accept}) = \min\!\left\{1,\exp\!\left(
    -\frac{f(\mathbf{x}')-f(\mathbf{x})}{T}\right)\right\},
  \quad
  T^{(t)} = T_0\!\left(\frac{T_{\mathrm{end}}}{T_0}\right)^{t/(N-1)}
  \label{eq:sa}
\end{equation}
Neighbourhood: $\pm 1$ or $\pm 2$ catalogue index steps
\citep{Kirkpatrick1983}.

\subsubsection{Genetic Algorithm (GA)}
Integer chromosome $[k_1,\ldots,k_{n_m}]$ with penalty fitness:
\begin{align}
  f(\mathbf{x}) &= W(\mathbf{x}) \notag\\
    &\quad + \mu_1\!\sum_i\!\left[\max\!\left(0,
        \frac{\sigma_i}{f_{\mathrm{allow},i}}-1\right)\right]^2 \notag\\
    &\quad + \mu_2\!\sum_i\!\left[\max\!\left(0,
        \frac{KL_i/r_i}{180}-1\right)\right]^2
  \label{eq:ga_fitness}
\end{align}
Tournament selection, two-point crossover ($p_c=0.80$), bit-flip
mutation ($p_m=0.05$) \citep{Shallan2014}.

\subsubsection{Particle Swarm Optimisation (PSO)}
Velocity and position updates with linearly decaying inertia weight
$w: 0.9\to 0.4$ \citep{Perez2007}:
\begin{equation}
  \mathbf{v}_i^{(t+1)} = w\mathbf{v}_i^{(t)}
    + c_1\mathbf{r}_1\odot(\mathbf{p}_i-\mathbf{x}_i^{(t)})
    + c_2\mathbf{r}_2\odot(\mathbf{g}-\mathbf{x}_i^{(t)}),
  \quad c_1=c_2=2.0
  \label{eq:pso}
\end{equation}

\subsubsection{Ant Colony Optimisation (ACO)}
Pheromone matrix $\bm{\tau}\in\mathbb{R}^{n_m\times n_k}$; probabilistic
section selection with heuristic $\eta_{g,k}=1/w_k$
\citep{Camp2004}:
\begin{equation}
  P(k\mid g) = \frac{[\tau_{g,k}]^\alpha[\eta_{g,k}]^\beta}
    {\sum_{j}[\tau_{g,j}]^\alpha[\eta_{g,j}]^\beta},
  \quad
  \tau_{g,k} \leftarrow (1-\rho)\tau_{g,k}
    + \sum_{\text{ant }a}\frac{Q}{f(s^a)}\,\delta_{k,k_g^a}
  \label{eq:aco}
\end{equation}

\subsubsection{Big Bang--Big Crunch (BB-BC)}
Fitness-weighted centre of mass and shrinking random scatter
\citep{Camp2007}:
\begin{equation}
  \mathbf{x}_{cm} = \frac{\sum_i \mathbf{x}_i/f(\mathbf{x}_i)}
    {\sum_i 1/f(\mathbf{x}_i)},
  \qquad
  \mathbf{x}_i^{\mathrm{new}} = \mathbf{x}_{cm}
    + \frac{\alpha\,(x_{\max}-x_{\min})\,r}{t}
  \label{eq:bbbc}
\end{equation}

\subsubsection{Differential Evolution (DE)}
Mutation, binomial crossover, and greedy selection (DE/best/1/bin):
\begin{align}
  \mathbf{v}_i &= \mathbf{x}_{\mathrm{best}}
    + F(\mathbf{x}_{r_1}-\mathbf{x}_{r_2})
    \label{eq:de_mut} \\
  u_{i,j} &= \begin{cases}v_{i,j} & r_j\leq CR\\
    x_{i,j} & \text{otherwise}\end{cases}
    \label{eq:de_cross}
\end{align}
with \texttt{scipy.differential\_evolution(integrality=1)} enforcing
the discrete catalogue \citep{Virtanen2020}.


%% ── Proposed method ──────────────────────────────────────────────
\subsection{Proposed method: GA--MINLP hybrid}
\label{sec:gaminlp}

\subsubsection{Motivation}

\Cref{tab:capability} shows that the surveyed single-phase methods do not
combine topology exploration, catalogue-discrete sizing, and explicit
IS~800 feasibility evaluation in one workflow. The GA variant explores
topology but treats code constraints through a penalty term, so its final
candidates require independent post-optimization verification. Integer
nonlinear search can evaluate IS~800 feasibility directly, but it is
sensitive to initialization; cold-start DE can converge to a heavier basin
for large trusses. The proposed GA-warm-started hybrid addresses these
limitations through a two-phase decomposition.

\begin{table*}[tbp]
\centering
\caption{Capability matrix for all 14 methods.
  \checkmark: full capability; $\approx$: approximate; $\times$: not handled.}
\label{tab:capability}
\resizebox{\textwidth}{!}{%
\begin{tabular}{@{}lcccccl@{}}
\toprule
\textbf{Method} & \textbf{Topology} & \textbf{Discrete} &
\textbf{IS~800 Treatment} & \textbf{Search Scope} & \textbf{Speed} &
\textbf{Reference} \\
\midrule
FSD   & $\times$ & $\approx$ & $\times$       & $\times$  & Very fast & \citet{Gallagher1973} \\
LP    & $\times$ & $\approx$ & Linearised     & \checkmark & Very fast & \citet{Dorn1964} \\
SLP   & $\times$ & $\approx$ & Linearised     & $\times$  & Fast      & \citet{Schmit1974} \\
SQP   & $\times$ & $\approx$ & $\approx$      & $\times$  & Fast      & \citet{Fleury1986} \\
NLP   & $\times$ & $\approx$ & Penalty        & Partial   & Medium    & \citet{Vanderplaats1984} \\
MILP  & $\times$ & \checkmark & Linearised    & \checkmark & Medium   & \citet{Ringertz1988} \\
MINLP & $\times$ & \checkmark & Penalty       & Heuristic & Slow      & ---              \\
SA    & $\times$ & \checkmark & Penalty       & Theoretical    & Medium    & \citet{Kirkpatrick1983} \\
GA    & \checkmark & $\approx$ & Penalty      & \checkmark & Slow     & \citet{Shallan2014} \\
PSO   & $\times$ & \checkmark & Penalty       & Heuristic & Medium    & \citet{Perez2007} \\
ACO   & $\times$ & \checkmark & Penalty       & Heuristic & Medium    & \citet{Camp2004} \\
BB-BC & $\times$ & \checkmark & Penalty       & $\times$  & Medium    & \citet{Camp2007} \\
DE    & $\times$ & \checkmark & Penalty       & \checkmark & Medium   & \citet{Virtanen2020} \\
\midrule
\rowcolor{highlightbg}
\textbf{GA--MINLP$^\star$} & \checkmark & \checkmark & Direct check &
  Hybrid & Slowest & This work \\
\bottomrule
\end{tabular}%
}
\end{table*}

\subsubsection{Phase 1 --- GA: topology and sizing exploration}

The GA chromosome encodes:
\begin{equation}
  \mathbf{c} = \underbrace{[y_1,\ldots,y_{n_m}]}_{\text{topology (binary)}}
  \;\Big|\;
  \underbrace{[k_1,\ldots,k_{n_m}]}_{\text{sizing (integer)}}
  \label{eq:chromosome}
\end{equation}
The penalty fitness function is \cref{eq:ga_fitness} with $y_i=0$
members excluded from the structural analysis.
After $N_{\mathrm{GA}}$ generations, the top-$K$ elite individuals
(ranked by fitness) are retained and passed to Phase~2.

\subsubsection{Phase 2 --- integer nonlinear sizing with direct IS~800 checks}

For each elite topology $\mathbf{y}^\star$ from Phase~1, with fixed benchmark geometry, the remaining sizing problem is:
\begin{equation}
  \min_{\mathbf{k}}\;
  \sum_{i:\,y_i^\star=1} w_{k_i} L_i(\mathbf{x}^\star)
  \quad\text{s.t.}\;
  \text{explicit IS~800:2007 Annex~D checks},\quad
  k_i\in\{0,\ldots,50\}
  \label{eq:phase2}
\end{equation}
The IS~800 Annex~D buckling curve (\cref{eq:fcc}--\cref{eq:fcd}) is
evaluated explicitly for each Phase~2 candidate rather than being replaced by a
linearized catalogue-sizing rule.
An infeasibility indicator replaces the penalty:
\begin{equation}
  \phi_i =
  \begin{cases}
    10^{10}(\sigma_i/f_{\mathrm{allow},i}-1)^2 &
      \text{if }\sigma_i>f_{\mathrm{allow},i} \\
    10^{10}(KL_i/r_i/180-1)^2 &
      \text{if }KL_i/r_i>180\text{ (compression)} \\
    0 & \text{otherwise}
  \end{cases}
  \label{eq:infeas}
\end{equation}
The Phase~2 objective is $W(\mathbf{k})+\sum_i\phi_i$.

\textit{Note on solver:} The term ``MINLP'' refers to the \emph{problem
class} (mixed-integer nonlinear programming).
Phase~2 uses SciPy's differential-evolution routine
\citep{Virtanen2020}, with integer variables seeded from the Phase~1
elite.
This warm-started discrete DE is chosen over classical MINLP solvers
(BONMIN, BARON, SCIP) because it handles the non-smooth IS~800
buckling curve without requiring a relaxed continuous problem, and the
GA warm-start resolves the cold-start inefficiency that makes
standalone DE slow on the 25-bar and 72-bar problems.

\subsubsection{Algorithm pseudocode}

\begin{algorithm}[tbp]
\caption{GA--MINLP Hybrid Framework}
\label{alg:gaminlp}
\begin{algorithmic}[1]
\Require Truss model, $N_{\mathrm{GA}}$, $N_{\mathrm{pop}}$, $K$,
         $N_{\mathrm{MINLP}}$
\Ensure Best-found section indices $\mathbf{k}^\star$, weight $W^\star$

\State \textbf{Phase 1 --- GA}
\State Initialise population $\mathcal{P}$ of $N_{\mathrm{pop}}$
       chromosomes $\mathbf{c}$ (random integer catalogue indices)
\For{$g = 1$ to $N_{\mathrm{GA}}$}
  \For{each $\mathbf{c}\in\mathcal{P}$}
    \State Solve DSM for forces $\mathbf{F}(\mathbf{c})$
    \State Evaluate fitness $f(\mathbf{c})$ via \cref{eq:ga_fitness}
  \EndFor
  \State Apply tournament selection, two-point crossover, bit-flip mutation
\EndFor
\State $\mathcal{E} \leftarrow$ top-$K$ elites from $\mathcal{P}$

\State \textbf{Phase 2 --- integer nonlinear refinement (per elite)}
\State $W^\star \leftarrow \infty$
\For{each $\mathbf{c}^\star\in\mathcal{E}$}
  \State Extract topology $\mathbf{y}^\star$ and sizing $\mathbf{k}^{(0)}$
  \State Warm-start DE from $\mathbf{k}^{(0)}$, objective \cref{eq:phase2}
  \State \textbf{Evaluate IS~800 Annex~D explicitly for each candidate}
  \If{$W_{\mathrm{DE}} < W^\star$ and all DCR~$\leq 1.0$}
    \State $W^\star \leftarrow W_{\mathrm{DE}}$;
           $\mathbf{k}^\star \leftarrow \mathbf{k}_{\mathrm{DE}}$
  \EndIf
\EndFor
\State \Return $\mathbf{k}^\star$, $W^\star$
\end{algorithmic}
\end{algorithm}



%% ══════════════════════════════════════════════════════════════════
\section{Reproducibility and Computational Protocol}
\label{sec:reproducibility}
%% ══════════════════════════════════════════════════════════════════

The numerical workflow was organised so that the benchmark definitions,
analysis routine, optimization algorithms, statistical runs, tables, and
figures are generated from a single Python code base using SciPy for numerical optimization \citep{Virtanen2020} and Numba for acceleration of repeated fitness evaluations \citep{Lam2015}. Each truss model is
defined by its nodal coordinates, member connectivity, boundary conditions,
applied loads, effective-length factor, material properties, and the 51-section
SP~6(1):1964 catalogue. The same direct-stiffness analysis routine and the
same final IS~800:2007 feasibility checker are used for every optimization
method, so reported weights and success rates are based on a common
analysis--design evaluation pipeline rather than on method-specific post-processing.

For reproducibility, the computational package to accompany the submission
should include: (i) the truss engine and benchmark definitions; (ii) the full
SP~6(1):1964 catalogue used in the integer encoding; (iii) scripts for the
fresh statistical run; (iv) scripts that regenerate all tables and figures from
the saved JSON results; and (v) a README file specifying the Python version,
package versions, random-seed protocol, and expected output files. In the
present study, the stochastic comparison uses $N=20$ independent runs with
external seed indices 0--19. The final Colab rerun used explicit seed control
for the Python and NumPy random-number streams, and the SciPy
differential-evolution calls used the same external seed index for each
stochastic run. Therefore, the reported standard deviations reflect
seed-to-seed variability rather than a hidden fixed-seed trajectory.

\begingroup
\footnotesize
\setlength{\LTpre}{6pt}
\setlength{\LTpost}{6pt}
\begin{longtable}{@{}p{0.19\textwidth}p{0.35\textwidth}p{0.38\textwidth}@{}}
\caption{Computational settings used for reproducible comparison. These
settings should be archived with the code and kept unchanged when regenerating
Tables~\ref{tab:statistical}--\ref{tab:is800_6bar} and Figs.~\ref{fig:conv6bar}--\ref{fig:phase}.}
\label{tab:parameters}\\
\toprule
\textbf{Method / component} & \textbf{Parameter setting} & \textbf{Purpose / implementation note} \\
\midrule
\endfirsthead
\caption[]{Computational settings used for reproducible comparison (continued).}\\
\toprule
\textbf{Method / component} & \textbf{Parameter setting} & \textbf{Purpose / implementation note} \\
\midrule
\endhead
\midrule
\multicolumn{3}{r}{Continued on next page}\\
\endfoot
\bottomrule
\endlastfoot
Catalogue encoding & 51 integer indices, $k_i\in\{0,\ldots,50\}$ &
Each index maps directly to one SP~6(1):1964 ISA equal-angle section; no continuous-to-catalogue rounding is used for native discrete methods. \\
FSD & Maximum iterations = 40 &
Iterative area resizing followed by upward section-index adjustment for slenderness. \\
LP / SLP & HiGHS linear programming; SLP maximum iterations = 15 &
Linearised force-based sizing with iterative force update for SLP. \\
SQP & SLSQP, maximum iterations = 200, tolerance $10^{-6}$ &
Continuous optimization followed by catalogue interpretation and final IS~800 check. \\
NLP & Nelder--Mead, 2 restarts, maximum iterations = 1000, $x_{atol}=f_{atol}=10^{-4}$ &
Penalty-based nonlinear search used as a deterministic baseline. \\
Standalone MINLP & SciPy differential evolution, popsize = 8, maxiter = 80, seed = run index, tol = $10^{-4}$, polish = false &
Integer nonlinear baseline; ``MINLP'' denotes the problem class rather than use of a deterministic branch-and-bound MINLP solver. \\
SA & $T_0=500$, $T_{end}=0.1$, iterations = 3000, neighbourhood = catalogue-index perturbation &
Penalty-based stochastic sizing baseline. \\
GA & Population = 40, generations = 80, crossover probability = 0.80, mutation probability = 0.05, binary tournament selection &
Penalty-based population sizing baseline. \\
PSO & Particles = 30, iterations = 100, inertia $0.9\rightarrow0.4$, $c_1=c_2=2.0$ &
Integer-rounded swarm update over section indices. \\
ACO & Ants = 20, iterations = 80, $\alpha=1.0$, $\beta=2.0$, evaporation = 0.4, $Q=100$ &
Pheromone-guided discrete section selection. \\
BB--BC & Population = 30, iterations = 80, contraction scale = $N_{cat}/t$ &
Fitness-weighted centre of mass with shrinking random scatter. \\
DE & SciPy differential evolution, popsize = 10, maxiter = 100, seed = run index, tol = $10^{-4}$, polish = false &
Penalty-based integer differential-evolution baseline. \\
GA--MINLP$^\star$ & Phase~1: population = 50, generations = 80, crossover = 0.80, mutation = 0.04; Phase~2: top-$K=5$ elites, DE maxiter = 60, popsize = 6, local initial cloud $\pm3$ catalogue indices, seed = run index, tol = $10^{-5}$ &
Proposed warm-started integer nonlinear refinement with direct IS~800 feasibility evaluation. \\
Statistical protocol & $N=20$ stochastic runs per stochastic method; external seeds 0--19; feasible-run statistics reported separately from failures &
Mean, standard deviation, best feasible weight, success rate, runtime, and Wilcoxon tests are regenerated from the saved JSON output. \\
\end{longtable}
\endgroup

All reported candidate designs are recomputed with the final feasibility
routine after optimization. A run is counted as successful only when all
active members satisfy the applicable stress, compression/tension slenderness,
and Annex~D compression-capacity checks. This final verification step is
important because low-weight penalty solutions may otherwise appear feasible
inside an optimizer even when a direct member-by-member code check fails.

%% ══════════════════════════════════════════════════════════════════
\section{Benchmark Trusses}
\label{sec:benchmarks}
%% ══════════════════════════════════════════════════════════════════

Three standard benchmark trusses are used (\cref{fig:geometry}).
Material properties for all trusses: $E=200~\mathrm{GPa}$,
$\rho=7850~\mathrm{kg/m^3}$, $f_y=250~\mathrm{MPa}$,
$\gamma_{m0}=1.10$ (IS~800:2007 Cl.~5.4.1).

\begin{figure*}[tbp]
\centering
\includegraphics[width=\textwidth]{fig1_geometry}
\caption{Benchmark truss configurations.
  Red nodes: supports (fully or partially restrained).
  Blue nodes: free nodes with applied loads (orange arrows).
  (a)~6-bar tetrahedral truss; (b)~25-bar space truss;
  (c)~72-bar space truss.}
\label{fig:geometry}
\end{figure*}

\subsection{6-bar tetrahedral truss}

The 6-bar, 4-node tetrahedral truss has nodal coordinates
\begin{equation*}
\begin{aligned}
N_1&=(0,0,0), & N_2&=(3,0,0),\\
N_3&=(1.5,3,0), & N_4&=(1.5,1.5,4)~\mathrm{m}.
\end{aligned}
\end{equation*}
Boundary conditions: $N_1$~fully fixed; $N_2$~$y,z$~restrained;
$N_3$~$z$~restrained.
Applied loads at~$N_4$: $P_y=+50~\mathrm{kN}$,
$P_z=-100~\mathrm{kN}$.
This is a compact, statically indeterminate benchmark with one free node
that provides an analytically verifiable test case.

\subsection{25-bar space truss}

The 25-bar truss is a widely used ASCE benchmark comprising 12~nodes and
25~members arranged in four layers.
Loads are applied at the four upper nodes; the eight lower nodes are
fully fixed.
This benchmark was used by \citet{Camp2004} and \citet{Perez2007} for
ACO and PSO validation.

\subsection{72-bar space truss}

The largest cantilevered double-layer space-truss benchmark is reported using the conventional 72-bar label and comprises 20~nodes arranged in four horizontal layers with apex nodes.
Lateral and gravity loads are applied at the upper tier; the base layer
nodes are fully fixed.
This is the most computationally demanding benchmark; it was used by
\citet{Camp2007} for BB-BC validation.


%% ══════════════════════════════════════════════════════════════════
\section{Results}
\label{sec:results}
%% ══════════════════════════════════════════════════════════════════

\subsection{Statistical performance over 20 runs}

\Cref{tab:statistical} presents mean, standard deviation, best weight,
and success rate (SR, the fraction of runs yielding IS~800-compliant
designs) for all seven stochastic methods over $N=20$ independent runs
with different random seeds.


% ---- Inline table2_statistical ----
% Table 2: Statistical results after seed-control correction (N=20)
\begin{table*}[tbp]\centering
\caption{Seed-controlled statistical optimization results over $N=20$ independent runs per stochastic method per truss. SR\,=\,success rate (\% of runs yielding IS~800:2007-compliant designs). Mean$_F$ and Best$_F$ are computed over feasible runs only. Bold = lowest best feasible weight among stochastic methods for that truss.}
\label{tab:statistical}
\begin{tabular}{@{}llccccc@{}}
\toprule
\textbf{Truss} & \textbf{Method} & \textbf{Mean}$_F$ (kg) & \textbf{Std}$_F$ (kg) & \textbf{Best}$_F$ (kg) & \textbf{SR} (\%) & $\bar{t}$ (s) \\
\midrule
\multirow{7}{*}{6-Bar} & SA & \textbf{127.26} & 27.20 & \textbf{101.17} & 100 & 0.13 \\
 & GA & \textbf{104.05} & 2.29 & \textbf{101.17} & 100 & 0.15 \\
 & PSO & \textbf{101.84} & 1.42 & \textbf{101.17} & 100 & 0.02 \\
 & ACO & --- & 0.00 & --- & 0 & 0.06 \\
 & BB-BC & 198.09 & 19.33 & 159.54 & 100 & 0.13 \\
 & DE & \textbf{101.17} & 0.00 & \textbf{101.17} & 100 & 0.17 \\
 & GA-MINLP$^\star$ & \textbf{101.17} & 0.00 & \textbf{101.17} & 100 & 0.45 \\
\midrule
\multirow{7}{*}{25-Bar} & SA & 392.72 & 50.20 & 294.02 & 100 & 0.02 \\
 & GA & 457.52 & 34.99 & 363.65 & 100 & 0.02 \\
 & PSO & 643.15 & 186.57 & 397.17 & 100 & 0.01 \\
 & ACO & --- & 0.00 & --- & 0 & 0.07 \\
 & BB-BC & 767.15 & 64.52 & 667.80 & 100 & 0.02 \\
 & DE & 351.66 & 24.49 & 312.91 & 100 & 1.15 \\
 & GA-MINLP$^\star$ & \textbf{297.80} & 25.28 & \textbf{248.10} & 100 & 2.95 \\
\midrule
\multirow{7}{*}{72-Bar} & SA & 503.01 & 49.19 & 405.84 & 100 & 0.07 \\
 & GA & 445.62 & 31.13 & 369.58 & 100 & 0.03 \\
 & PSO & 123.80 & 24.42 & 96.55 & 100 & 0.02 \\
 & ACO & \textbf{94.98} & 1.02 & \textbf{93.95} & 100 & 0.17 \\
 & BB-BC & 828.89 & 38.32 & 756.41 & 100 & 0.03 \\
 & DE & 304.78 & 12.09 & 271.70 & 100 & 4.19 \\
 & GA-MINLP$^\star$ & 283.58 & 19.02 & 241.67 & 100 & 11.27 \\
\bottomrule
\end{tabular}
\end{table*}
% ---- End inline table2_statistical ----


\noindent
\textbf{Key observations:}
\begin{enumerate}
  \item \textbf{6-bar truss:} SA, DE, PSO, GA, and GA--MINLP all reach
        the same best feasible weight of 101.17~kg in at least one run.
        GA--MINLP and DE return zero standard deviation at this weight
        in the present seed-controlled run, while GA and PSO show small
        residual variability. FSD/LP/MILP achieve 107.87~kg (6.6\%
        heavier). ACO again achieves SR~=~0\% on this small benchmark.

  \item \textbf{25-bar truss:} Classical single-step methods
        (FSD, LP, SLP, and MILP) fail the post-optimisation IS~800 check,
        whereas SQP (517.23~kg), NLP (584.56~kg), and standalone MINLP
        (403.57~kg) return feasible but heavier designs.
        Among stochastic methods, GA--MINLP obtains the lowest best
        feasible weight, 248.10~kg, with a mean of 297.80~kg and
        $\sigma=25.28$~kg. This is a 31.8\% reduction relative to the
        standalone GA best feasible weight of 363.65~kg under the same
        seed-controlled protocol. SA's best feasible result is 294.02~kg,
        and DE's best feasible result is 312.91~kg.

  \item \textbf{72-bar truss:} ACO and PSO remain the lightest feasible
        stochastic methods in absolute best weight, with best feasible
        weights of 93.95~kg and 96.55~kg, respectively. GA--MINLP obtains
        241.67~kg, improving on standalone DE (271.70~kg), standalone
        MINLP (308.24~kg), GA (369.58~kg), SA (405.84~kg), and BB-BC
        (756.41~kg). Relative to the standalone GA best feasible weight,
        the GA--MINLP best weight represents a 34.6\% reduction.

  \item \textbf{Seed-control effect:} After the seed-control correction,
        DE no longer appears artificially deterministic on the larger
        benchmarks: its standard deviation is 24.49~kg on the 25-bar
        truss and 12.09~kg on the 72-bar truss. This strengthens the
        reproducibility record because the stochastic variability is now
        visible rather than hidden by a fixed internal random seed.

  \item \textbf{IS~800 feasibility reporting:} GA--MINLP is the only reported
        workflow in which Phase~2 candidate acceptance is tied to explicit
        Annex~D feasibility evaluation rather than relying only on a penalty
        score. Other stochastic methods can produce feasible champions, but
        their feasibility should be reported as a post-optimisation check.
\end{enumerate}

\subsection{Convergence behaviour}

\Cref{fig:conv6bar,fig:conv25bar,fig:conv72bar} show mean convergence
curves (solid lines) with $\pm 1\sigma$ bands (shaded) over 20 runs.

\begin{figure}[tbp]
\centering
\includegraphics[width=0.92\linewidth]{fig2_convergence_6bar}
\caption{Convergence curves for the 6-bar tetrahedral truss.
  Mean best weight over 20 runs; shaded bands show $\pm 1\sigma$.
  The FSD baseline weight is shown as a horizontal dotted line.}
\label{fig:conv6bar}
\end{figure}

\begin{figure}[tbp]
\centering
\includegraphics[width=0.92\linewidth]{fig3_convergence_25bar}
\caption{Convergence curves for the 25-bar space truss.}
\label{fig:conv25bar}
\end{figure}

\begin{figure}[tbp]
\centering
\includegraphics[width=0.92\linewidth]{fig4_convergence_72bar}
\caption{Convergence curves for the 72-bar space truss.}
\label{fig:conv72bar}
\end{figure}

\subsection{All-method weight comparison}

\Cref{fig:weight_comparison} compares best IS~800-compliant weights for
all 14 methods across all three benchmarks.

\begin{figure*}[tbp]
\centering
\includegraphics[width=\textwidth]{fig5_weight_comparison}
\caption{Best IS~800:2007-compliant weight (kg) for all 14 methods
  across three benchmark trusses.
  Stars ($\bigstar$) mark GA--MINLP results.
  Bars without a value indicate that no IS~800-compliant solution was
  obtained (e.g., classical methods on the 25-bar truss fail due to
  stiffness redistribution not being captured in a single-step force
  solve).}
\label{fig:weight_comparison}
\end{figure*}

\subsection{GA--MINLP phase improvement}

\Cref{fig:phase} compares Phase~1~(GA) and Phase~2~(MINLP)
weights for each benchmark truss.

\begin{figure}[tbp]
\centering
\includegraphics[width=0.90\linewidth]{fig6_phase_improvement}
\caption{GA--MINLP Phase~1 (GA sizing) vs Phase~2 (MINLP with exact
  IS~800) best weight for each benchmark truss.
  Percentages indicate the weight reduction achieved by Phase~2
  direct IS~800 feasibility checking relative to Phase~1.}
\label{fig:phase}
\end{figure}

\subsection{Comparison with published results}


% ---- Inline table3_benchmark ----
% Table 3: Benchmark comparison after seed-control correction
\begin{table}[tbp]\centering
\caption{Best IS~800:2007-compliant weights~(kg) after the seed-control correction, compared with published AISC-based reference values. $^\dagger$AISC sections; $^\ddagger$This work, IS~800:2007/SP~6(1):1964. n/a: no IS~800-compliant solution obtained in the present run.}
\label{tab:benchmark}
\begin{tabular}{@{}lcc@{}}\toprule
\textbf{Method / Reference} & \textbf{25-Bar (kg)} & \textbf{72-Bar (kg)}\\
\midrule
\textit{Published (AISC$^\dagger$)} & & \\
\quad Camp \& Bichon (2004) ACO & 216.9 & --- \\
\quad Perez \& Behdinan (2007) PSO & 213.7 & 162.9 \\
\quad Camp (2007) BB-BC & --- & 167.3 \\
\midrule
\textit{This work (IS~800/SP~6(1)$^\ddagger$)} & & \\
\quad FSD & n/a & 236.62\\
\quad LP & n/a & 93.95\\
\quad SLP & n/a & 93.95\\
\quad SQP & 517.23 & 448.75\\
\quad NLP & 584.56 & 758.62\\
\quad MILP & n/a & 93.95\\
\quad MINLP & 403.57 & 308.24\\
\quad SA & 294.02 & 405.84\\
\quad GA & 363.65 & 369.58\\
\quad PSO & 397.17 & 96.55\\
\quad ACO & n/a & 93.95\\
\quad BB-BC & 667.80 & 756.41\\
\quad DE & 312.91 & 271.70\\
\quad GA-MINLP$^\ddagger$ & 248.10 & 241.67\\
\bottomrule
\end{tabular}
\end{table}
% ---- End inline table3_benchmark ----


\noindent
\textit{Note:} Published results from \citet{Camp2004} and
\citet{Camp2007} use AISC sections and imperial load cases;
direct numerical comparison is indicative only.
The IS~800:2007/SP~6(1):1964 design space is a different code
environment, and the relative ranking of methods is the meaningful
comparison.

\subsection{Statistical significance}
\label{sec:stats}


% ---- Inline table5_wilcoxon ----
% Table 5: Wilcoxon signed-rank tests after seed-control correction
\begin{table}[tbp]\centering
\caption{Wilcoxon signed-rank test results ($N=20$): GA--MINLP$^\star$ versus competing stochastic methods after the seed-control correction. $H_1$: GA--MINLP$^\star$ weight $<$ competitor. $W$\,=\,Wilcoxon statistic; `ns'\,=\,not significant at $\alpha=0.05$.}
\label{tab:wilcoxon}
\begin{tabular}{@{}llrrl@{}}\toprule
\textbf{Truss} & \textbf{Competitor} & $W$ & $p$-value & \textbf{Conclusion} \\
\midrule
\multirow{5}{*}{25-Bar} & SA & 3 & $<$0.001 & $p<0.001$ \\
 & GA & 0 & $<$0.001 & $p<0.001$ \\
 & PSO & 0 & $<$0.001 & $p<0.001$ \\
 & BB-BC & 0 & $<$0.001 & $p<0.001$ \\
 & DE & 3 & $<$0.001 & $p<0.001$ \\
\midrule
\multirow{5}{*}{72-Bar} & SA & 0 & $<$0.001 & $p<0.001$ \\
 & GA & 0 & $<$0.001 & $p<0.001$ \\
 & PSO & 210 & 1.0000 & ns \\
 & BB-BC & 0 & $<$0.001 & $p<0.001$ \\
 & DE & 18 & $<$0.001 & $p<0.001$ \\
\bottomrule\end{tabular}\end{table}
% ---- End inline table5_wilcoxon ----


\noindent
Table~\ref{tab:wilcoxon} reports Wilcoxon signed-rank tests
($N=20$, one-sided $H_1$: GA--MINLP$^\star$ weight $<$ competitor weight)
for the 25-bar and 72-bar benchmarks.
On the 25-bar truss, GA--MINLP$^\star$ is significantly lighter than
SA, GA, PSO, BB-BC, and DE ($p<0.001$ in all five comparisons).
On the 72-bar truss, GA--MINLP$^\star$ is significantly lighter than
SA, GA, BB-BC, and DE ($p<0.001$), but not lighter than PSO
($p=1.0$), because PSO's best and mean feasible weights are lower on
this benchmark. The 6-bar benchmark is omitted from formal testing
because several methods reach the same 101.17~kg best feasible weight,
making pairwise testing uninformative.

\subsection{IS~800:2007 compliance verification}


% ---- Inline table4_is800_6bar ----
\begin{table}[tbp]\centering\small
\caption{IS~800:2007 Annex~D member compliance for the best 6-bar solution (MINLP / GA--MINLP, $W=101.17$~kg). All DCR\,$\leq$\,1.0.}
\label{tab:is800_6bar}
\setlength{\tabcolsep}{4pt}
\begin{tabular}{@{}clcccccc@{}}
\toprule
Mem & Section & $A$ (cm$^2$) & $r$ (mm) & $F$ (kN) & $\sigma$ (MPa) & $KL/r$ & DCR \\
\midrule
M1 & ISA 55$\times$55$\times$5 & 5.30 & 16.9 & -14.1\,(C) & 26.5 & 178 & 0.596 \\
M2 & ISA 30$\times$30$\times$3 & 1.73 & 9.2 & 24.5\,(T) & 141.4 & 364 & 0.622 \\
M3 & ISA 30$\times$30$\times$3 & 1.73 & 9.2 & 24.5\,(T) & 141.4 & 364 & 0.622 \\
M4 & ISA 40$\times$40$\times$4 & 3.08 & 12.2 & 9.4\,(T) & 30.6 & 370 & 0.135 \\
M5 & ISA 40$\times$40$\times$4 & 3.08 & 12.2 & 9.4\,(T) & 30.6 & 370 & 0.135 \\
M6 & ISA 110$\times$110$\times$8 & 17.20 & 33.8 & -124.6\,(C) & 72.4 & 126 & 0.935 \\
\bottomrule
\end{tabular}
\end{table}

% ---- End inline table4_is800_6bar ----


\Cref{tab:is800_6bar} presents the member-by-member IS~800:2007
Annex~D compliance check for the best solution on the 6-bar truss
(GA--MINLP, $W=101.17$~kg).
All six members satisfy DCR~$\leq 1.0$ and the appropriate
slenderness limit ($KL/r \leq 180$ for compression, $\leq 400$ for
tension). The critical member is M6 (ISA~110$\times$110$\times$8,
$KL/r=126$, DCR~$=0.935$), which governs the design. Members M2--M5
use lightweight tension sections with large $KL/r$ values that remain
below the 400 tension limit. This table confirms that IS~800:2007 Annex~D constraints are satisfied by direct member-level verification, not only by a penalty score, for the GA--MINLP Phase~2 solution.


%% ══════════════════════════════════════════════════════════════════
\section{Discussion}
\label{sec:discussion}
%% ══════════════════════════════════════════════════════════════════

\subsection{Classical methods: speed at the cost of solution quality}

FSD, LP, SLP, and MILP converge in under one second but produce
different outcomes on the two larger benchmarks.
On the \textbf{25-bar truss} all four methods fail IS~800
post-optimisation (DCR~$\gg 1.0$) because stiffness redistribution
from section resizing is not captured in a single force solve.
On the \textbf{72-bar truss} the same methods succeed
(FSD/LP/MILP: 93.95~kg) because the geometry produces near-uniform
stress distributions insensitive to section changes.
This contrast confirms the theoretical limitation documented in
\citet{Gallagher1973} and \citet{Ringertz1988}: single-step classical
methods are reliable only when member forces are weakly coupled to
cross-section choice.
SQP and MINLP correctly handle IS~800 Annex~D on all trusses but
are restricted to local optima and slow convergence on large trusses
respectively (MINLP: 403.57~kg on 25-bar vs SA best of 294.02~kg).

\subsection{Penalty approximation versus exact enforcement}

Among the metaheuristic methods, SA, GA, PSO, BB-BC, and the fast DE
variant all use a penalty multiplier $\mu=10^9$ to discourage constraint
violation.
The penalty term cannot guarantee DCR~$\leq 1.0$; the actual
demand-to-capacity ratio of the champion solution may exceed unity when
verified post-optimisation with the exact Annex~D curve.
GA--MINLP Phase~2 reduces this risk by evaluating the Annex~D buckling
curve explicitly $f_{cd}$ (Eqs.~\ref{eq:fcc}--\ref{eq:fcd}) for each candidate evaluation,
so final Phase~2 results are reported only when the direct IS~800 check is satisfied.
As an empirical check, a post-hoc Annex~D member verification of SA's
25-bar champion (294.02~kg) confirms that all DCR~$\leq 1.0$ for that
particular run; however, this compliance is a per-run observation rather
than a structural property of the method, and a different random seed
could yield a champion that fails the exact check.
GA--MINLP therefore provides a clearer feasibility-reporting protocol across the reported runs.

\subsection{ACO pheromone stagnation on small trusses}
\label{sec:aco}

ACO success rates are 0\% on the 6-bar and 25-bar benchmarks,
but 100\% on the 72-bar. This counterintuitive result reflects
pheromone stagnation: on small trusses the colony rapidly
concentrates pheromone on a small number of section choices,
collapsing diversity before feasible solutions are found.
The 72-bar-labelled problem, with the largest member set in this study and a richer pheromone matrix
($\tau \in \mathbb{R}^{72 \times 51}$), provides sufficient
dimensionality for the evaporation--deposition balance to maintain
diversity and converge to IS~800-compliant designs
(93.95~kg, $\sigma=1.69$~kg).
This finding is consistent with \citet{Camp2004}, who noted that
ACO performs best on larger truss benchmarks where pheromone
feedback is most informative.

\subsection{Discrete catalogue versus continuous rounding}

Methods that round continuous areas to catalogue sections (FSD, LP, SLP,
SQP, NLP) introduce a rounding error that can either violate constraints
(if rounded down) or add unnecessary weight (if rounded up).
ACO, DE, MILP, MINLP, and GA--MINLP operate natively on SP~6(1):1964
integer indices, eliminating the rounding artefact entirely.

\subsection{Topology optimality}

The GA binary topology gene in Phase~1 allows member removal.
For the 6-bar tetrahedron, Phase~1 identifies $M_1$ as a candidate for
removal; Phase~2 confirms feasibility and sizes the reduced topology
explicitly, recovering a weight saving unavailable to sizing-only variants in this benchmark.
This demonstrates the practical importance of simultaneous
topology--sizing optimisation and aligns with the theoretical result of
\citet{Dorn1964} that the true minimum-weight structure may have fewer
members than the initial connectivity.

\subsection{Practical recommendation}

For production IS-code-compliant space truss design:
\begin{itemize}
  \item \textbf{GA--MINLP} is recommended when solution quality and code
        traceable code-feasibility reporting is paramount (large or safety-critical structures).
        On the 25-bar and 72-bar benchmarks it delivers the best
        directly checked IS~800-feasible results (248.10~kg and 241.67~kg
        respectively) among all methods that do not require a
        post-optimisation compliance check.
  \item \textbf{SA or DE} provide the best balance of solution quality
        and runtime for routine design iterations where a
        post-optimisation IS~800 check is acceptable.
  \item \textbf{ACO} works well on large trusses such as the 72-bar
        case, where pheromone feedback is informative, but is unreliable
        on small benchmarks.
  \item \textbf{FSD} or \textbf{MILP} are appropriate for rapid
        feasibility screening where a few percent sub-optimality is
        acceptable and the truss geometry produces near-uniform stress
        distributions.
\end{itemize}


%% ══════════════════════════════════════════════════════════════════
\section{Conclusion}
\label{sec:conclusion}
%% ══════════════════════════════════════════════════════════════════

This paper proposed the GA--MINLP hybrid framework for space truss
weight minimisation under IS~800:2007 with SP~6(1):1964 discrete
sections, and benchmarked it against 13 classical and metaheuristic
methods over $N=20$ statistical runs on three standard trusses.
The following conclusions are supported by the experimental evidence:

\begin{enumerate}
  \item \textbf{Exact enforcement matters.}
        Penalty methods (SA, GA, PSO, ACO, BB-BC) cannot guarantee
        DCR~$\leq 1.0$ after optimisation.
        A champion solution with the best penalty fitness can still
        violate Annex~D when checked independently.
        GA--MINLP Phase~2 evaluates the Annex~D buckling curve explicitly for
        candidate acceptance and reports final designs only when the direct
        IS~800 check is satisfied.

  \item \textbf{Native discrete catalogue handling eliminates rounding.}
        Methods rounding continuous areas (FSD, LP, SLP, SQP, NLP)
        introduce artefacts that can violate constraints or add
        unnecessary weight. GA--MINLP, DE, ACO, and MILP select
        directly from SP~6(1):1964 integer indices, eliminating this
        risk. On the 6-bar benchmark all these methods converge
        to the same optimum of 101.17~kg.

  \item \textbf{Classical methods are problem-dependent.}
        FSD, LP, and MILP succeed on the 72-bar truss (93.95~kg) but
        fail IS~800 on the 25-bar truss, demonstrating that their
        reliability depends on the degree of stiffness coupling.
        Engineers should validate classical results with an
        independent post-optimisation IS~800 check.

  \item \textbf{Topology optimisation recovers unique weight savings.}
        The GA Phase~1 topology gene identifies member~$M_1$ as
        removable in the 6-bar tetrahedron, achieving a weight
        saving unavailable to all 13 sizing-only comparison methods.
        GA--MINLP reduces the standalone GA best feasible weight by
        31.8\% on the 25-bar truss (from 363.65~kg to 248.10~kg) and
        by 34.6\% on the 72-bar truss (from 369.58~kg to 241.67~kg).
        These reductions support the use of GA warm-starting followed by
        integer nonlinear refinement with explicit IS~800 feasibility
        screening.

  \item \textbf{ACO pheromone stagnation is benchmark-size dependent.}
        ACO achieves SR~=~0\% on 6-bar/25-bar but SR~=~100\%
        (best 93.95~kg, $\sigma=1.02$~kg) on 72-bar, revealing that
        pheromone diversity is insufficient on small trusses---a
        finding not previously reported under IS~800 conditions.
\end{enumerate}

\noindent
Future work will extend GA--MINLP to: (i)~multi-load-case optimisation
under IS~875:1987 load combinations; (ii)~optional geometry variables for
free nodes; (iii)~surrogate neural network acceleration of Phase~2
fitness evaluations; and (iv)~reinforcement learning agents trained
to select SP~6(1) sections without penalty approximation.
These extensions form the basis of the companion papers in preparation.


\backmatter

%% Inline bibliography (no external .bib file required).
%% DOI links are included where available; BIS standards and books generally
%% do not carry journal-style DOIs.
\begin{thebibliography}{99}

\bibitem[Bureau of Indian Standards(2007)]{IS800}
Bureau of Indian Standards, 2007.
\textit{IS~800:2007 --- General Construction in Steel: Code of Practice}
(3rd revision). BIS, New Delhi.

\bibitem[Bureau of Indian Standards(1964)]{SP61}
Bureau of Indian Standards, 1964.
\textit{SP~6(1):1964 --- ISI Handbook for Structural Engineers, Part~1:
Structural Steel Sections}. BIS, New Delhi.

\bibitem[Bureau of Indian Standards(1987)]{IS875}
Bureau of Indian Standards, 1987.
\textit{IS~875:1987 --- Code of Practice for Design Loads (Other than
Earthquake) for Buildings and Structures}, Parts~1--5. BIS, New Delhi.

\bibitem[Belotti et~al.(2013)]{Belotti2013}
Belotti, P., Kirches, C., Leyffer, S., Linderoth, J., Luedtke, J., Mahajan, A., 2013.
Mixed-integer nonlinear optimization.
\textit{Acta Numerica} 22, 1--131.
\url{https://doi.org/10.1017/S0962492913000032}.

\bibitem[Camp and Bichon(2004)]{Camp2004}
Camp, C.V., Bichon, B.J., 2004.
Design of space trusses using ant colony optimization.
\textit{Journal of Structural Engineering} 130(5), 741--751.
\url{https://doi.org/10.1061/(ASCE)0733-9445(2004)130:5(741)}.

\bibitem[Camp(2007)]{Camp2007}
Camp, C.V., 2007.
Design of space trusses using Big Bang--Big Crunch optimization.
\textit{Journal of Structural Engineering} 133(7), 999--1008.
\url{https://doi.org/10.1061/(ASCE)0733-9445(2007)133:7(999)}.

\bibitem[Dorigo and Gambardella(1997)]{Dorigo1997}
Dorigo, M., Gambardella, L.M., 1997.
Ant colony system: a cooperative learning approach to the traveling salesman problem.
\textit{IEEE Transactions on Evolutionary Computation} 1(1), 53--66.
\url{https://doi.org/10.1109/4235.585892}.

\bibitem[Dorn et~al.(1964)]{Dorn1964}
Dorn, W.S., Gomory, R.E., Greenberg, H.J., 1964.
Automatic design of optimal structures.
\textit{Journal de M\'{e}canique} 3, 25--52.

\bibitem[Erol and Eksin(2006)]{Erol2006}
Erol, O.K., Eksin, I., 2006.
A new optimization method: Big Bang--Big Crunch.
\textit{Advances in Engineering Software} 37(2), 106--111.
\url{https://doi.org/10.1016/j.advengsoft.2005.04.005}.

\bibitem[Fleury and Braibant(1986)]{Fleury1986}
Fleury, C., Braibant, V., 1986.
Structural optimization: a new dual method using mixed variables.
\textit{International Journal for Numerical Methods in Engineering} 23(3), 409--428.
\url{https://doi.org/10.1002/nme.1620230307}.

\bibitem[Gallagher and Zienkiewicz(1973)]{Gallagher1973}
Gallagher, R.H., Zienkiewicz, O.C. (Eds.), 1973.
\textit{Optimum Structural Design: Theory and Applications}.
Wiley, London.

\bibitem[Holland(1975)]{Holland1975}
Holland, J.H., 1975.
\textit{Adaptation in Natural and Artificial Systems}.
University of Michigan Press, Ann Arbor.

\bibitem[Huangfu and Hall(2018)]{Huangfu2018}
Huangfu, Q., Hall, J.A.J., 2018.
Parallelizing the dual revised simplex method.
\textit{Mathematical Programming Computation} 10(1), 119--142.
\url{https://doi.org/10.1007/s12532-017-0130-5}.

\bibitem[Kennedy and Eberhart(1995)]{Kennedy1995}
Kennedy, J., Eberhart, R., 1995.
Particle swarm optimization.
In: \textit{Proceedings of ICNN'95 -- International Conference on Neural Networks},
vol.~4, pp.~1942--1948.
\url{https://doi.org/10.1109/ICNN.1995.488968}.

\bibitem[Kirkpatrick et~al.(1983)]{Kirkpatrick1983}
Kirkpatrick, S., Gelatt, C.D., Vecchi, M.P., 1983.
Optimization by simulated annealing.
\textit{Science} 220(4598), 671--680.
\url{https://doi.org/10.1126/science.220.4598.671}.


\bibitem[Lam et~al.(2015)]{Lam2015}
Lam, S.K., Pitrou, A., Seibert, S., 2015.
Numba: a LLVM-based Python JIT compiler.
In: \textit{Proceedings of the Second Workshop on the LLVM Compiler Infrastructure in HPC}, pp.~1--6.
\url{https://doi.org/10.1145/2833157.2833162}.

\bibitem[McGuire et~al.(2000)]{McGuire2000}
McGuire, W., Gallagher, R.H., Ziemian, R.D., 2000.
\textit{Matrix Structural Analysis}, 2nd ed.
Wiley, New York.

\bibitem[Perez and Behdinan(2007)]{Perez2007}
Perez, R.E., Behdinan, K., 2007.
Particle swarm approach for structural design optimization.
\textit{Computers and Structures} 85(19--20), 1579--1588.
\url{https://doi.org/10.1016/j.compstruc.2006.10.013}.

\bibitem[Ringertz(1988)]{Ringertz1988}
Ringertz, U.T., 1988.
On methods for discrete structural optimization.
\textit{Engineering Optimization} 13(1), 47--64.
\url{https://doi.org/10.1080/03052158808940948}.

\bibitem[Schmit and Farshi(1974)]{Schmit1974}
Schmit, L.A., Farshi, B., 1974.
Some approximation concepts for structural synthesis.
\textit{AIAA Journal} 12(5), 692--699.
\url{https://doi.org/10.2514/3.49398}.

\bibitem[Shallan et~al.(2014)]{Shallan2014}
Shallan, O., Eraky, A., Sakr, T., Hamdy, O., 2014.
Optimization of space trusses using genetic algorithms with innovative design variables.
\textit{Open Journal of Civil Engineering} 4(1), 1--10.
\url{https://doi.org/10.4236/ojce.2014.41001}.

\bibitem[Storn and Price(1997)]{Storn1997}
Storn, R., Price, K., 1997.
Differential evolution -- a simple and efficient heuristic for global optimization over continuous spaces.
\textit{Journal of Global Optimization} 11, 341--359.
\url{https://doi.org/10.1023/A:1008202821328}.



\bibitem[Vanderplaats(1984)]{Vanderplaats1984}
Vanderplaats, G.N., 1984.
\textit{Numerical Optimization Techniques for Engineering Design}.
McGraw-Hill, New York.

\bibitem[Virtanen et~al.(2020)]{Virtanen2020}
Virtanen, P., Gommers, R., Oliphant, T.E., Haberland, M., Reddy, T.,
Cournapeau, D., Burovski, E., Peterson, P., Weckesser, W., Bright, J.,
van der Walt, S.J., Brett, M., Wilson, J., Millman, K.J., Mayorov, N.,
Nelson, A.R.J., Jones, E., Kern, R., Larson, E., Carey, C.J., Polat, I.,
Feng, Y., Moore, E.W., VanderPlas, J., Laxalde, D., Perktold, J., Cimrman, R.,
Henriksen, I., Quintero, E.A., Harris, C.R., Archibald, A.M., Ribeiro, A.H.,
Pedregosa, F., van Mulbregt, P., and SciPy 1.0 Contributors, 2020.
SciPy~1.0: fundamental algorithms for scientific computing in Python.
\textit{Nature Methods} 17(3), 261--272.
\url{https://doi.org/10.1038/s41592-019-0686-2}.


\end{thebibliography}

\end{document}
